\documentclass{article}
\usepackage{graphicx}
\usepackage{subcaption}
\begin{document}

\begin{figure}[h]
\centering
  \begin{subfigure}[b]{0.45\textwidth}
    \centering
    \fbox{\rule{0pt}{5cm}\rule{5cm}{0pt}}  % placeholder (replace with \includegraphics)
    \caption{Subfigure 1 caption}
    \label{fig:sub1}
  \end{subfigure}
  \hfill
  \begin{subfigure}[b]{0.45\textwidth}
    \centering
    \fbox{\rule{0pt}{5cm}\rule{5cm}{0pt}}
    \caption{Subfigure 2 caption}
    \label{fig:sub2}
  \end{subfigure}

  \vspace{1em}

  \begin{subfigure}[b]{0.45\textwidth}
    \centering
    \fbox{\rule{0pt}{5cm}\rule{5cm}{0pt}}
    \caption{Subfigure 3 caption}
    \label{fig:sub3}
  \end{subfigure}
  \hfill
  \begin{subfigure}[b]{0.45\textwidth}
    \centering
    \fbox{\rule{0pt}{5cm}\rule{5cm}{0pt}}
    \caption{Subfigure 4 caption}
    \label{fig:sub4}
  \end{subfigure}

\caption{This is a figure containing several subfigures.}
\label{fig:subfigs}
\end{figure}

In the text, you can refer to subfigures of Figure~\ref{fig:subfigs}
as \ref{fig:sub1}, \ref{fig:sub2}, \ref{fig:sub3}, and \ref{fig:sub4},
and to the sub-index as (a), (b), (c) and (d).
\end{document}